\section{Spiritual LARPing}

\begin{quotex}
The others, the mages, the theosophists, the cabalists, the spiritists, the hermetics, the Rosicrucians, remind me, when they are not mere thieves, of children playing and scuffling in a cellar. And if one descends lower yet, into the hole-in-the-wall places of the pythonesses, clairvoyants, and mediums, what does one find except agencies of prostitution and gambling? All these pretended peddlers of the future are extremely nasty; that's the only thing in the occult of which one can be sure. \flright{\textsc{J K Huysmans}, \emph{Là-Bas}}

\end{quotex}
\paragraph{Role Playing}
There is the mistaken impression that esoteric studies are based on some secret knowledge, hidden from others. \textbf{Valentin Tomberg}, however, makes a distinction between a secret and an arcanum.

\begin{quotex}
Secrets are only facts, procedures, practices, or whatever doctrines that one keeps to oneself for a personal motive, since they are able to be understood and put into practice by others to whom one does not want to reveal them. 

\end{quotex}
You see, it has nothing to do with esoteric attainment; rather, it is a means to retain control and power. Those who are in on the secret can use it to identify in and out groups. At its worst, it only feeds the urge for self-inflation. For these reasons, we prefer to keep everything in the open.

It is not the text or the book that ultimately matters. The text is merely the reflection of the pure act of intelligence, as revealed in the Magician Arcanum. It is only possible to reach the pure act through spiritual exercises. Otherwise, as Tomberg points out, you are only doing philosophy at best, or just role playing at worst.

Our first goal here is to promote the ideal of self-mastery and to show ways to achieve it. This requires control of thinking, feeling, and willing. This control can only result from a purification of one's thinking, feeling, and willing. Without such purification, the mirror will be perturbed, and its reflections will distort the pure act of intelligence.

I have known and spoken to, in my travels, gurus, lamas, anthroposophists, theosophists, freemasons, rosicrucians, priests, gurdjieffians, Jungians, charismatics, channelers, and so on. Even giving the benefit of the doubt, many of the followers are role playing, i.e., they enjoy the fantasy of spiritual attainment without having to do the work. They try to act as they think a spiritual person should act. They dress a certain way, adopt a new diet, say the right words, while leaving their inner life untouched.

In actuality, you would not readily identify an initiate even if he crossed your path. He would have no reason to reveal himself to you, unless you expressed a sincere interest in the teachings.

\paragraph{The Christian Hermetic Tradition}
Our task has been to serve the living Tradition, which has been assigned by Valentin Tomberg to his unknown friends. To that end, we have described the various manifestations of that Tradition throughout time and space. We have also introduced a great many authors who have contributed to that Tradition. The goal has been to entice others to want to serve that same Tradition. We have been quite open with that material, since it hides itself from those unable or unwilling to understand it.

As a matter of principle, it is problematic to earn one's living from Hermetic teaching. It can lead to subtle distortions based on the need ``give the customer what he wants". At its lowest, it is a form of simony, i.e., selling spiritual knowledge for a price. In the olden days, a Hermetist might take the role of a street performer or horse trader. In that way, he could travel from village to village without arising suspicions, and meet with other Hermetists. It would also provide a modest income. That was a motivation for the Tarot cards. Books were too expensive and bulky to travel with. So a deck of cards that encapsulated the teachings in symbolic form was a boon.

Once upon a time, someone was hectoring us every now and then to rewrite or even remove certain posts, to delete comments, or to block links. It made no sense, until I was informed of the reason:

\begin{quotex}
I thought it necessary to free Tomberg from rightist view and people who follow you 

\end{quotex}
That is absurd. No one has the authority to decide who can read this blog or comment on Tomberg. First of all, we don't have any political or religious tests to join our groups. Instead, we rely on the Law of Affinity; that will bring the right people together if they belong together.

The Traditional way is not partisan, so the distinction between ``right" and ``left" is a meaningless creation of modernity. Tradition is the path of unity whereas partisanship is dual, hence devilish. \textbf{Louis Claude de Saint-Martin} explains:

\begin{quotex}
Now, in order to show how [the numbers] are related to their base of activity, let us begin by observing the working of unity and of the number two. When we contemplate an important truth, such as the universal power of the Creator, his majesty, his love, his profound light, or suchlike attributes, we bear ourselves wholly towards this supreme model of all things; all our faculties are suspended in order to fill us with him, and we really only make ourselves one with him. This is the active image of unity, and the number one in our languages is the expression of this unity or invisible union which, existing intimately between all attributes of this unity, must equally exist between it and all its produced creations. But if, after having borne all our faculties of contemplation towards this universal source, we return our gaze to ourselves and fill ourselves with our own contemplation, in such a way that we regard ourselves as the origin of some of the inner light or satisfaction that this source has procured for us, from that moment we establish two centres of contemplation, two separate and rival principles, two bases which ate not linked; lastly, we establish two unities, with this difference—that one is real and the other is apparent … 

But to divide being through the middle is to divide it into two parts; it is to pass from the whole to the quality of the part or the half, and it is here that the true origin of illegitimate twofoldness lies. . .this example is sufficient to show us the birth of the number two — to show us the origin of evil. . . 

\end{quotex}
\paragraph{Choosing a Path}
In a recent comment, Mikkel asked a very good question:

\begin{quotex}
How does one truly understand what ``I want"? Due to the story of my experiences of my life, I find that I am more blind to this unfortunately in periods of distress. 

\end{quotex}
Only by coming to that realisation, is it possible for a man to seek a way out. Most men live in ignorance of their real nature and what is truly good for them.

The Parable of the Coach explains our situation: our body is inefficient, our emotions are uncontrolled, our intellect is asleep, and there is no one in charge. Hence the coach moves without purpose or goal, in random directions.

It may happen that a more awake coachman may pass by and offer help. Usually the help is refused. Nevertheless, some few, due perhaps to karma, fortune, or grace, are willing to accept such help.

It sounds simple, except that in our world there are many offers to help, most asking for money. So it is right to be sceptical. People may spend more effort in choosing a car mechanic or doctor than in their spiritual development. It is reasonable to know what results to expect. The answer is an improvement in one's thinking, feeling, and willing functions. You should ask those who have been helped previously about their experiences. The results should be permanent, not just some emotional high after a rousing lecture.

Avoid cults. Apply this test when in doubt:

\begin{quotex}
A cult is easy to join, but hard to leave. A free group is hard to find, but easy to leave. 

\end{quotex}
If the goal is self-mastery, then clearly you cannot surrender your will to someone else. You need to understand your situation and learn to choose your own goals. Simply being told what to do is useless. Either you will dislike the advice and leave the teaching, or else you will like the advice and not bother with the exercises.

A contrary approach is offered by many others. Typically, it will include spiritual counselling, for a fee. The counsellor will then tell the student his karmic tasks and what his life goals should be. To repeat, the student needs to determine that on his own, even if it might require a little encouragement. I have never seen that approach work very well.

\paragraph{Uniting Hermetists}
I've always hoped for a loose coalition of Hermetists, each developing the Hermetic Tradition in its own way, despite the attacks described above. These different groups would cooperate in various ways, without regard for border. Tomberg's description is what I have in mind:

\begin{quotex}
Spiritual exercises in common form the common link that unites Hermetists. It is not knowledge in common which unites them, but rather the spiritual exercises and the experience which goes hand in hand with them. If three people from different countries were to meet each other, having made the book of Genesis by Moses, the Gospel of St. John, and the vision of Ezekiel, the subject of spiritual exercises for many years, they would do so in brotherhood. … What one knows is the result of personal experience and orientation, whereas depth, the level one attains — disregarding the aspect and extent of knowledge that one has gained — is what one has in common. 

\end{quotex}
Tomberg tells us how to evaluate a teaching:

\begin{quotex}
There are mirages above, as there are mirages below; you only know that which is verified by the agreement of all forms of experience in its totality—experience of the senses, moral experience, psychic experience, the collective experience of other seekers for the truth, and finally the experience of those whose knowing merits the title of wisdom and whose striving has been crowned by the title of saint. 

\end{quotex}
In other words, blind belief based solely on authority or on unverifiable claims are not helpful on the Hermetic path. Too much emphasis on unusual or fantastic experiences will discourage students who should be content with slow but steady progress. Steiner, in \emph{Knowledge of Higher Worlds}, warns several times about the expectation of fantastical experiences, for example he writes:

\begin{quotex}
it is imperative to extirpate the idea that any fantastic, mysterious practices are required for the attainment of higher knowledge. It must be clearly realized that a start has to be made with the thoughts and feelings with which we continually live, and that these feelings and thoughts must merely be given a new direction. 

\end{quotex}
Moreover, it is best to observe the Hermetic principle, ``Be Silent". Boasting of spiritual attainments, even if they are true, is not helpful. Steiner, in the same book, verifies that thought:

\begin{quotex}
know how to observe silence concerning your spiritual experiences. 

\end{quotex}
\paragraph{Mysticism and Initiation}
\textbf{Rene Guenon} distinguishes between mysticism and initiation. The former is passive, individualistic, subjective, and exoteric. Initiation, on the other hand, is quite different. He explains that

\begin{quotex}
Initiation is the transmission, through appropriate rites, of a spiritual influence 

\end{quotex}
\textbf{Titus Burkhardt} expands on the idea that the transmission of the spiritual influence must be conferred

\begin{quotex}
by a master who also communicates the method and confers the means of spiritual concentration that are appropriate to the aptitudes of the disciple 

\end{quotex}
Furthermore, he points out that in ancient times a mystic referred to someone who had knowledge of the mysteries. In other words, the mystic was initiated in the esoteric mysteries, but that is not how the word is used in our time.

\paragraph{Preconditions for Spiritual Science}
A period of self-purification is necessary before any effective initiation can take place. Otherwise, all sorts of mischief and misunderstandings can result when the unprepared receive esoteric teachings. First of all, it is best to keep one's physical body in good order, lest it bring distractions to spiritual practice. Proper diet and exercise are helpful. Moreover, the concentration required for physical disciplines like athletics or martial arts are a good preparation.

One's life should also be in order. Anxiety over relationships or money will make spiritual practice difficult. Then, work can be done on oneself. The Christian Hermetist \textbf{Willi Seiss}, in his commentary on Valentin Tomberg's \emph{Course on the Lord's Prayer} claims:

\begin{quotex}
A precondition is the purification of the three soul-forces of willing, feeling, and thinking. 

\end{quotex}
He summarizes the specifics like this:

\begin{itemize}
\item The element of will is trained through the rigorous, many-sided process of cleansing from everything that is impure. 
\item Through refraining from destructive feelings, that is, all reprehensible emotions 
\item Disciplining of the thoughts, whereby the essential is freed of influences from fearful, corrupt, hate-filled and empty thoughts. 
\end{itemize}
\paragraph{The Pagan Revival}
In our time, the neo-pagans are prime examples of this. Misunderstanding the cycles and subcycles in the flow of time, as well as the existence of castes, they claim to hearken back to the exoterism of their allegedly pagan ancestors. Anyone else, they say, are bugmen. Yet that have no idea what caste they descend from. Moreover, since they believe they are the ``true" Traditionalists, they do not bother with the idea of purification, and don't really understand the inner sense of the pagans. The pagans also had their bugmen.

It is true enough, however, that the strata of earlier pagan traditions remain in our psyches. The \textbf{Book of Revelation} contains messages about the ancient Indians, Persian, Egyptian, Greek, and Roman cultures. Their psychic residues remain in our collective unconscious. They must both be remembered and transcended.

\paragraph{Note}
Friday the 13th was the day the \textbf{Templars }were destroyed. This post is dedicated to them.

\emph{And God knows best.}



\flrightit{Posted on 2019-09-13 by Cologero }

\begin{center}* * *\end{center}

\begin{footnotesize}\begin{sffamily}



\texttt{James on 2018-09-23 at 10:54 said: }

Your illness comes up in my meditative moments every now and then . When I focused on it , it always seemed to center around the idea of ``filtration" . Kidneys , obviously , serve to filter the subtle substance of our blood .

Whenever I would ponder on it , I thought about how this illness affects that organ which is meant to filter and clean the liquid that is so often associated with the ``esoteric" current in man . It seemed ``appropriate" considering how so many self-proclaimed esotericists seem to only be ``poisoning" the stream and there is no organ capable of filtering it out for the body of humanity .

The interesting thing about the Dante reading that I do is when I think about the circle of Violence and how in one part of it — those who blaspheme — are filled with so much internal burning that they burn more than the flames that descend from the ``sky" . I can't help but imagine that condemning someone to Hell is some kind of blasphemy (the irony is not lost on me considering Dante allegorically does ``exactly this" though the thin line is clear: Dante is conversing in the language of an ``open secret" where he ``condemns" people to Hell almost as a ``test" for us to look beyond the particular person at the eternal work that he's created and what is being written about is not some canonization of individuals to Hell , but a canonization of sin into Hell — something which an eye looking at only phenomenon and not essences cannot achieve when reading Dante . This is part of the reason why Dante's greyhound and a mangy , emaciated she-wolf (the symbols of truth and deception respectively) look so ``similar" [they are both so `thin’] . Dante is speaking of one of the central mysteries of Christianity: ``why is what is good so often so indistinguishable from what is bad superficially" [e.g. the symbol of death and torture being the symbol of salvation ; God dying being the vehicle for salvation . etc]) . Still , Dante was quite instructive to me in how the sin of blasphemy manifests: that it creates the hellish burning inside the person as they persist in that attitude . Thus , such a statement as , ``you will die , die in horror , die for ages , you will burn in hell like no other" is already indicative of a burning interior life .

Naturally , I'm not claiming some kind of moral superiority nor am I gawking at a sinful act considering how much of a hypocritical sinner I can be , but I was merely pausing to think of it and realizing just how contemporary Dante's trek through Hell can be for me . Here , we pass yet another sinner suffering the effects of the sin itself . The beauty of Free Will is that it's so easy to see the Hell we create if we but look with the Pilgrim's eyes . The elegant beauty is that we truly get not just what we deserve but also what we wanted . Hell is the oddly poetic and beautiful synthesis of both free choice and justice . In eternity , we never escape these singularities whether they flow upward or downward . Everything always reaches synthesis . As Flannery O'Connor had put it as her title: ``everything that rises must converge" . Obviously the opposite is also true : ``everything that descends must converge" (tying it back to Dante's Satan) . The Author of Meditations speaks of the same thing with the synthesis to white and the chromatic loss below when he used a colour analogy .

Still , I would consider what your friend spoke about to be interesting . She must consider you to be the blasphemer (I'm only guessing with limited information) . I wonder if , perhaps , she imagines that you do not treat the subject matter with enough rituality . That's not to say she's necessarily wrong even if her reaction might be counterproductive to her aims (again , I'm only guessing off grossly incomplete assumptions , but I decided to follow it to some `logical' conclusion just to see if it might be at all helpful) .

In my encounters with Jungians , for example , there have been quite a few methods proposed to me that ought to be done ``properly" and plenty have given me their prescriptions on what this ``propriety" means . To do otherwise would invite danger and collapse . A friend of mine recently said that an encounter with one's Shadow or Anima , for example , could produce years of psychosis (and he might be right) . Nonetheless , I'm honestly interested in why your friend decided that you had been guilty of some kind of sacrilege that you deserve condemnation . It is not that I intend to take sides (that is rather puerile anyway) , but , rather , I wish to ``understand" .

As for gornahoor itself . I find it fascinating . Perhaps it would have acted as a filter (to tie it back to the kidneys) to try and find one or two or three or however many in order to contribute to the ``communal soul of art , science , and religion" together , but I personally , and perhaps quite wrongly , always saw it attract a certain type of individual who was more interested in using esoteric ``knowledge" as clubs to swing at ``others" . Evola become the new Nietschze for fiesty young people unable to accept reality .

I think your meditation on death and this whole endeavour with Gornahoor is instructive . After all , what I find most interesting about the appeal to sectarianism is that it often exposes a very clear insecurity: the fear of death .

‘"We" [whoever the ``we" happen to be] ``are afraid to perish from this earth" .'

I have said it once before a long time ago : perhaps the solution is not to live by the sword and die by it , but to learn to accept one's cross in order to be resurrected . Esoteric knowledge seems to be used as a weapon by many who frequent these halls in order to scream by at the encroaching darkness . Perhaps these individuals need to learn how to properly accept death . It is not a call to cowardice or complacency , but a call to finding that strength in inner peace rather than inner fear .


\hfill

\texttt{Mikkel on 2018-09-23 at 14:40 said: }

``my advice is to do what you want to do now, and not later."

How does one truly understand what ``I want"? Do to the story of my experiences of my life, I find that I am more blind to this unfortunately in periods of distress. I think this comes in time with self-mastery…perhaps Evola (are we still speaking of him now?) said it well in the translation posted years ago (cited above):

``Against psychoanalysis the ideal of a Self (or ``I") must be validated, a Self who does not abdicate, who intends to remain conscious, autonomous, and sovereign in the confrontation with the dark and subterranean part of his soul and the demon of sensuality; that does not mean either ``repression" or a psychotic scission, but he brings about an equilibrium of all his faculties ordered to a higher meaning of living and acting."

Many individuals seem to want to ``dominate" or be in power over those dark parts of his soul, but is that truly equilibrium? The analogy of the chariot/coach is far more apt, there are drivers and what drives (horses), but in totality it is a living system that must function together. Usually understanding and respect in a relationship goes much farther than the whip hand.

``There are some Christian Hermetic groups that prefer to work in private. In our opinion, since Tomberg that is no longer appropriate." What changed after Tomberg for this to be true? No true line to carry on his work? Sometimes others are far more open than we give them credit to a surprising extent or maybe I've been careful of who to associated myself with, the Spirit breatheth where he will.

Other times while reading about betrayals like the one mentioned above I notice the melancholy that surfaces within. I look forward to reading the material that comes out from the translation, my German is quite basic so I am sorry that I cannot be of much help. Similarly, I would love to take all of those books, but the sense of responsibility of research with those feels heavy in comparison to my current talents. Knowledge is more similar to that concept than I realized…we look to it earnestly at times, but what happens to us when we do nothing with it? Are we not to help others in some way? Maybe I need to learn my own place, isn't it always like that when one comes down to the essence of understanding any problem?

You and those of the `community' are in my thoughts and prayers, this I know is true and something I want to do.


\hfill

\texttt{Cologero on 2018-09-23 at 16:07 said: }

Mikkel, Tomberg made clear that the unknown friends are those who dare to plunge into the currents of Christian Hermetism. There is no secret group that controls who can be initiated or determines what ``grade" someone is at. Any such organisation that tries to control a teaching always ends up protecting itself not the teaching. The Hermetist does not offer an alternative to religion nor to science, rather he strives to integrate all knowledge into a larger whole.


\hfill

\texttt{Cologero on 2018-09-23 at 16:27 said: }

James, perhaps I am a blasphemer in someone's eyes. Of course, even a blasphemer deserves to know the charges and offer a defense.

Or maybe it is a difference in approach. Again, you put me in a position that may seem boastful, but there have been holy fools, zen masters, etc., whose techniques may not meet certain standards of decorum. Then again, a man shouldn't take himself too seriously in these matters.

Maybe you should add Nassim Taleb to your reading list. Knowledge is not usually gained from the top down, but rather from a lot of experience. The Romans built aqueducts, sometime with a drop of 12 inches over the length of a mile. You probably know how to compute the slope, but I don't know how the Romans did that calculation with their numerical system. Nonetheless, water got to the city. Practicality trumps theory in this case.

That is why our method uses heuristics. We try a lot of ways to come to awareness; over time, one learns what works for him. All the theory in the world won't is not helpful, and is usually a distraction.

For example, suppose I tell you that by awakening a certain chakra or invigorating some of its petals, then you will achieve inner peace. Does that knowledge bring you inner peace? Of course not. First of all, it is necessary to achieve inner peace — I prefer the term ataraxia — and then perhaps you can see how the chakras change in response.

The person of deep faith, with no esoteric knowledge, can achieve such a state. That is my model.


\hfill

\texttt{Auld Wat on 2022-05-11 at 08:12 said: }

``First of all, it is best to keep one's physical body in good order, lest it bring distractions to spiritual practice. Proper diet and exercise are helpful. Moreover, the concentration required for physical disciplines like athletics or martial arts are a good preparation. One's life should also be in order. Anxiety over relationships or money will make spiritual practice difficult. Then, work can be done on oneself."

Thank you for the recommendations.


\end{sffamily}\end{footnotesize}
